
We applied LoRA to Mamba-130m exactly as we would to any transformer --- correct target modules, correct parameter counts, 1.14\% trainable weights --- and watched SST-2 accuracy sit at 50.9\% for three consecutive training epochs, indistinguishable from the zero-shot baseline of 49.1\%. This is not slow convergence. This is zero learning.

The result is surprising precisely because everything visible looked correct. The PEFT library installed the adapters without complaint. The model's state dictionary contained LoRA weight matrices at every targeted projection layer. The training loop ran, loss values were logged, and gradients were computed. And yet, after three full epochs over 4,000 training samples --- a protocol sufficient for LoRA to converge on SST-2 with transformers --- the model remained locked at majority-class prediction accuracy, as if fine-tuning had never occurred. The same pattern held across all four GLUE tasks evaluated: SST-2 (flat), MNLI (degrading), QNLI (oscillating near baseline), and QQP (consistently 0.000).

\subsubsection{1.1 The Surface Problem}

Parameter-efficient fine-tuning, and LoRA in particular \citep{Hu2022}, has become the dominant strategy for adapting large pretrained language models to downstream tasks. By freezing the base model and training only low-rank decompositions of selected weight matrices, LoRA achieves near-full-fine-tuning accuracy at a fraction of the trainable parameter cost. The method's appeal is partly its apparent architecture-agnosticism: any layer implemented as \texttt{nn.Linear} is a candidate for LoRA adaptation.

Mamba \citep{Gu2023} is a leading transformer-alternative for long-context language modeling. Its selective state space mechanism achieves linear-time training and constant-memory inference, making it a natural target for deployment scenarios where transformer attention is prohibitively expensive. Mamba-130m is publicly available on the HuggingFace Hub, its projection layers --- \texttt{in\textbackslash \_proj}, \texttt{out\textbackslash \_proj}, and \texttt{x\textbackslash \_proj} --- are all \texttt{nn.Linear}, and community implementations (alxndrTL/mamba-peft) report SST-2 accuracy of 90--92\% with LoRA rank 8 on these exact layers.

The natural inference is that LoRA on Mamba works. Practitioners building Mamba-based NLP pipelines would reasonably apply LoRA to these projection layers and expect learning to occur.

\subsubsection{1.2 The Deeper Problem and the Gap}

The assumption that LoRA on \texttt{nn.Linear} produces learning is valid for transformers because it rests on an implicit guarantee: the gradient path from the task loss to the LoRA weight matrices is unobstructed. In transformers, the attention mechanism is implemented as fully differentiable operations under standard PyTorch autograd. Any \texttt{nn.Linear} target in the attention stack receives well-defined gradients when trained on any loss function.

In Mamba, this guarantee does not hold in the same way. The architecture's core computation --- the selective state space scan \texttt{h\textbackslash \_t = \$Ā\textbackslash cdot h\_\{t-1\}\$ + B̄\$\textbackslash cdot x\_t\$} --- is implemented either as a custom CUDA parallel associative scan kernel (when the Mamba CUDA extensions are installed) or as a sequential Python fallback (when they are not). Our experiments ran under the sequential fallback implementation, as indicated by the log warning: ''The fast path is not available because one of \texttt{(selective\textbackslash \_state\textbackslash \_update, selective\textbackslash \_scan\textbackslash \_fn, causal\textbackslash \_conv1d\textbackslash \_fn, causal\textbackslash \_conv1d\textbackslash \_update, mamba\textbackslash \_inner\textbackslash \_fn)} is None. Falling back to the sequential implementation of Mamba.'' This implementation uses standard PyTorch operations, meaning gradient flow through the SSM computation is mediated by autograd --- yet the learning failure persists.

Our experiments demonstrate that this gradient path is ineffective for classification under the base checkpoint with a randomly-initialized classification head. Three training epochs produce exactly zero improvement over zero-shot accuracy on SST-2 (50.9\% vs. 49.1\% zero-shot, a net gain of 1.84 percentage points --- within statistical noise). MNLI accuracy actively degrades from its zero-shot level of 34.6\% to 32.3\% by the final epoch. The loss oscillates between approximately 0.65 and 0.73 across training batches without monotonic decrease --- the signature of incoherent gradient updates rather than directed learning.

No prior work, to our knowledge, provides a controlled, reproducible characterization of projection-only LoRA failure on pure Mamba SSMs for classification tasks. The community reference (alxndrTL/mamba-peft, 2024) reports 90--92\% SST-2 without fully documenting the training recipe --- which checkpoint (base vs. instruction-tuned), which classification head design, which prompt format, how many epochs. Practitioners attempting to replicate this result with the base \texttt{mamba-130m-hf} checkpoint, a standard classification head, and a clean training loop will observe what we observed: no learning. Without a clear failure characterization, the field will repeatedly rediscover this failure independently.

\subsubsection{1.3 Key Insight}

We observe that projection-only LoRA on Mamba-130m produces a characteristic failure pattern --- loss oscillation without monotonic decrease, combined with flat or degrading per-epoch accuracy --- that is consistent across all four GLUE tasks evaluated. The failure occurs even under the sequential (autograd-compatible) implementation of the SSM scan, suggesting that the gradient barrier is not exclusively a property of the CUDA kernel's backward pass, but may reflect the geometric properties of the gradient signal passing through a pretrained SSM scan applied to a randomly-initialized classification head.

The SSM scan was designed and optimized for next-token prediction, the task on which Mamba-130m was pretrained. For a randomly-initialized classification head placed on the final hidden state, the gradient is novel and apparently fails to translate into discriminative updates to the LoRA weight matrices, regardless of whether the scan is implemented in CUDA or PyTorch.

This is a conceptual distinction the PEFT community needs: API compatibility --- the ability to install LoRA adapters with correct parameter counts --- is a necessary but not sufficient condition for fine-tuning to produce learning. Gradient compatibility, verified empirically through per-epoch task accuracy, is the authoritative criterion.

\subsubsection{1.4 Contributions}

Our investigation reveals three contributions:

First, we provide direct empirical evidence that projection-only LoRA (targeting \texttt{in\textbackslash \_proj}, \texttt{out\textbackslash \_proj}, \texttt{x\textbackslash \_proj} with rank 8) does not produce meaningful classification task learning on Mamba-130m across all four evaluated GLUE tasks. The evidence is precise: SST-2 accuracy of 0.5092 across three identical consecutive training epochs (zero-shot baseline: 0.4908), MNLI accuracy degrading from 0.3463 zero-shot to 0.3234 across three epochs, QNLI oscillating near baseline (zero-shot: 0.5046; epochs 1/2/3: 0.5011/0.5149/0.5126), and QQP remaining at 0.0000 across all epochs. The overall GLUE average improves only marginally from 0.3354 (zero-shot) to 0.3363 (LoRA).

Second, we document verified zero-shot GLUE baselines for Mamba-130m: SST-2 = 0.4908, MNLI = 0.3463, QNLI = 0.5046, QQP = 0.0000. These baselines, confirmed across two independent evaluations in the experimental environment, provide a reproducible reference for future Mamba adaptation experiments.

Third, we characterize a diagnostic signature for practitioners --- loss oscillation without monotonic decrease, combined with flat or degrading per-epoch accuracy --- and note that this failure manifests under the sequential Python fallback implementation of the SSM scan (i.e., under standard PyTorch autograd), which has implications for how the mechanistic cause should be understood.

\subsubsection{1.5 Paper Organization}

Section 2 reviews Mamba's architectural properties, existing PEFT methods, and the MambaPEFT literature. Section 3 describes the MUST\_WORK existence check experimental design. Section 4 presents the experimental setup. Section 5 reports results across all four GLUE tasks. Section 6 interprets the results, addresses the sequential fallback finding, and discusses limitations. Section 7 concludes with directions for future work.

